This paper analyses a form of happiness that resists the standard theories of well-being: the plain, shared happiness in which nothing happens, nothing is achieved, and nothing is lacked, yet which is full rather than empty. Such happiness is neither a \emph{state} attained nor an \emph{event} undergone. The paper argues that it is instead an \emph{emergence}: a fullness arising from the way a generative relation continues itself, irreducible to its substrate, and whose form is fixed not in advance but historically and materially. A survey of the major traditions---joint-attention psychology, psychoanalysis, phenomenology, eudaimonism, neuroscience, structuralism, political economy, and process philosophy---shows that each reads such happiness either as a \emph{lack} or as a \emph{static satisfaction}, and so misses its distinctive feature: that this stillness is not stagnation but continued generation.

To articulate this, the paper develops a relational ontology in which the individual is a knot, not a premise, of relation, and locates the engine of generation, by way of psychoanalysis, in the irreducible gap between the symbolic and the real. The same gap can be lived in two modes: a mode of lack, in which the subject is driven to fill it, and a mode of fullness, in which the subject dwells in the openness it sustains. Happiness is the second mode. The paper further argues that happiness is \emph{hermeneutic}---given only as meaning to a subject of the symbolic---and therefore that no purely dynamical account can reach it; it demonstrates this by pressing the dynamical-systems vocabulary to its limit and showing it cannot adjudicate competing intuitions of happiness. The dynamical is nonetheless retained as the \emph{background manifold} on which meaning is laid down: meaning is not reducible to it, but cannot stand apart from it, so happiness is neither physical fact nor arbitrary reading.

From this follow two results. First, shared happiness is not the merger of two subjects into one symbolic order---hermeneutic annexation---but the resonance of two meaning-worlds that remain mutually translatable without being exhausted. Second, the silence proper to such happiness completes a series: the silence of obligation, of restraint, and of needlessness. Throughout, no single form is permitted to occupy the place of happiness's essence.
